\begin{tablalibro}{2}{L{0.31}L{0.69}}
\caption{Glosario general}\label{tab:glosario-general}\\
\toprule
\rowcolor{Papel}
\CabeceraTabla{Término} & \CabeceraTabla{Definición} \\
\midrule
\endfirsthead
\SigueTabla{2}
\toprule
\rowcolor{Papel}
\CabeceraTabla{Término} & \CabeceraTabla{Definición} \\
\midrule
\endhead
\midrule
\endfoot
\bottomrule
\multicolumn{2}{@{}p{\linewidth}@{}}{\NotaLibro{Elaboración propia.}}\\
\endlastfoot

Abstracción & Selección de características pertinentes de un problema. \\
Acumulador & Variable que conserva un resultado parcial. \\
Algoritmo & Procedimiento definido para resolver un problema. \\
Caso de prueba & Entrada y condiciones con resultado esperado. \\
Colección & Estructura que agrupa varios datos. \\
Complejidad & Crecimiento del costo computacional. \\
Condición & Expresión evaluada como verdadera o falsa. \\
Contador & Variable que registra ocurrencias. \\
Dato & Representación de información susceptible de procesamiento. \\
Depuración & Proceso de localizar y corregir defectos. \\
Descomposición & División de un problema en componentes. \\
Diagrama de flujo & Representación gráfica de un procedimiento. \\
Entrada & Información suministrada a una solución. \\
Función & Unidad de código invocable. \\
Generalización & Adaptación de una solución a casos de una misma clase. \\
Iteración & Repetición de un conjunto de instrucciones. \\
Patrón & Regularidad reconocible entre situaciones. \\
Pensamiento computacional & Enfoque de formulación y resolución sistemática de problemas. \\
Procesamiento & Transformación de entradas en resultados. \\
Pseudocódigo & Representación textual estructurada de un algoritmo. \\
Refactorización & Reorganización interna sin cambio observable de comportamiento. \\
Requisito & Condición o capacidad que debe satisfacer una solución. \\
Salida & Resultado generado por el procesamiento. \\
Validación & Comprobación de adecuación a necesidades. \\
Variable & Nombre asociado a un valor en un programa. \\
Verificación & Comprobación de cumplimiento de especificaciones. \\
\end{tablalibro}
